% !TEX root = A0.MiTFM.tex

\pagestyle{fancy}
\renewcommand{\headrulewidth}{0pt}
\chapter*{Abstract}
\addcontentsline{toc}{chapter}{Abstract}

Cloud computing is now a core component of modern technology, and the need to apply adequate security measures has become one of its most relevant challenges. Kubernetes is one of the main tools used to deploy cloud-native systems, but its relative complexity can make secure configuration difficult for all parties involved.

This work configures an intentionally vulnerable Kubernetes environment in order to simulate compromise scenarios on the host running the applications. The study focuses on the offensive exploitation of these weaknesses by using Kyverno, a Kubernetes policy engine, to perform attacks aimed at achieving persistence in the system. Existing policies are modified and new policies are created from an offensive security perspective. The attack planning is based on a STRIDE risk analysis intended to identify and exploit relevant vectors within the deployed environment.

\mbox{} \bigskip

\noindent \textbf{Keywords}:
\begin{compactitem}
    \item Cybersecurity.
    \item Cloud computing.
    \item Kubernetes.
    \item Kyverno.
    \item STRIDE.
\end{compactitem}

\blankpage
